\documentclass[../DoAn.tex]{subfiles}
\begin{document}

\begin{center}
    \Large{\textbf{ABSTRACT}}\\
\end{center}
\vspace{1cm}

This graduation thesis focuses on researching, designing, and building \textbf{JQK Study} - a comprehensive online learning and multiple-choice testing platform based on a modern Microservices architecture. In the context of digital transformation in education, online evaluation has become increasingly popular, but existing platforms face significant challenges: difficulty in monitoring student cheating during remote exams, time-consuming manual exam generation, and the fragmentation between learning management systems (LMS) and independent testing tools.

To address these challenges, this project proposes an integrated solution that unifies learning and testing workflows into a single platform. The system is implemented using a Microservices architecture, featuring high-performance \textbf{Go} for the backend services and \textbf{Next.js} with \textbf{TailwindCSS} for the frontend interface. The highlights of the proposed solution include a real-time, multi-layered anti-cheat monitoring mechanism. At the client-side, the platform captures events such as tab switching, window blurring, and copy/paste actions utilizing Page Visibility and Clipboard APIs to log violation attempts. At the server-side, the system maintains an independent countdown timer to prevent local browser tampering. Additionally, the system features an automated exam generation algorithm that randomly selects questions based on difficulty levels and topic structures defined by the instructor.

The final outcome of this thesis is a fully functional platform deployed on a containerized environment (Docker and Kubernetes) with core services including API Gateway, User Service, Course Service, Exam Service, and Notification Service. The system helps instructors save significant time in preparing exams, manage classroom progress, and provides visualized violation logs, ensuring fairness and integrity in online assessments.

\end{document}